\section{One source argument on two analytic domains}
\label{cp:foundations}

The following interface supplies the analytic approximation domains and
real-time bounds used by the compression constructions.  Its proof contains
the stopped-cavity argument and all its auxiliary coefficient estimates.

\begin{proposition}[Analytic sources and real-time response bounds]
\label{cp:source}
Fix the finite query lists below before initialization. Eventually at each
fixed confidence, the following hold on one event. In this proposition put
\[
 T=32\frac m\gamma\log(en),\qquad
 r_t=\frac{\gamma}{\beta^{30L}Y^2m\sqrt{(d+3)\log(en)}},\qquad
 r_q=\frac1{\beta^{3L}\sqrt{(d+3)\log(en)}}.
\]
Write \(k^{(L)}=w\), \(k^{(j)}=W^{(j+1)\top}\delta^{(j+1)}\) (\(j<L\)) for
the backward responses before the activation derivative. The four source
families are
\begin{equation}
 h^{(j)},\qquad W^{(j)}(0)h^{(j-1)},\qquad
 \delta^{(j)},\qquad W^{(j+1)}(0)^\top\delta^{(j+1)},
 \label{cp:source-families}
\end{equation}
omitting the forward image for \(j=1\) and reverse image for \(j=L\).

\begin{enumerate}[label=(\roman*)]
\item The forward and backward fields are holomorphic on a neighborhood of
\[
 ([-r_t,T+r_t]+i[-r_t,r_t])\times
 \{v\in\mathbb C^d:v^\top v=1,
                      \|\operatorname{Im}v\|_2\le\sinh r_q\}.
\]
Query holomorphy is relative to the quadric; for \(d=1\) use \(\{-1,1\}\).
Here and for \(t\in[0,\infty]\) over the real sphere, the first-layer
operator divided by \(\sqrt n\) and all hidden operators have norm at most
ten, and
\begin{equation}
 \frac{\|h^{(j)}\|_2}{\sqrt n}\le\beta^{3L},\qquad
 \frac{\max\{\|\delta^{(j)}\|_2,\|k^{(j)}\|_2\}}{\sqrt n}
 \le16\beta^{4L}\frac{Ym}{\gamma}.
 \label{cp:source-rms}
\end{equation}
Each source coordinate has modulus at most \(\beta^{6L}\sqrt n\)
on the complex domain.

\item On each fixed finite query list, the predictions are holomorphic near
\([-r,T+r]+i[-r,r]\), where
\begin{equation}
 r=\frac m{\gamma\sqrt{\log(en)}},\qquad
 |f_n(t,\sqrt d\,v)|\le16\beta^{6L}\frac{Ym}{\gamma}.
 \label{cp:finite-time}
\end{equation}
On the entire real training trajectory,
\begin{equation}
 \sup_{t\in[0,\infty]}\max_{a,j,i}|k_{a,i}^{(j)}(t)|
 \le32\beta^{21L}\frac{Ym}{\gamma}\sqrt{\log(en)}.
 \label{cp:carrier}
\end{equation}

\item For a declared finite panel containing the training inputs, there are
\(0=t_0<\cdots<t_J=T\) and \(r_b>0\), \(0\le b<J\), with
\(t_{b+1}-t_b\le r_b/2\). The sources are holomorphic near
\(|t-t_b|\le r_b\), with the operator, RMS and coordinate bounds in (i):
forward queries range over the panel, backward queries over training inputs.
The partition and radii depend deterministically only on
\(n,m,\gamma,Y\), activations and depth, and
\begin{equation}
 J\le1+\beta^{30L}\left(\frac{Ym}{\gamma}\right)^2\sqrt{\log(en)}
       +\beta^{16L}\frac m\gamma
             \log\left(1+\frac{4\log(en)}{\log2}\right).
 \label{cp:panel-source}
\end{equation}
\end{enumerate}
\end{proposition}

\begin{proof}
We prove stronger intermediate domains with explicitly calculated bounds,
then compare them with the displayed interface.  All symbols introduced
from here through the end of this proof are private to this argument;
the subsequent constructions use only the proposition's stated bounds.
All probability statements concern one sufficiently large width, with the
data, depth, activations and confidence fixed.  Write
\(\lambda=\gamma/m\), \(Y=\|y\|_2/\sqrt m\),
\(S=16Y/\lambda\), and \(\ell=\log(en)\).  We treat \(Y>0\);
zero labels give the stationary zero predictor.  For each full or cavity
trajectory write \(\rho=\|r\|_2/\sqrt m\), also at complex times using
the Euclidean norm of the complex residual vector.  Write
\(W_0^{(j)}=W^{(j)}(0)\) for initialized matrices.  Put
\[
 b=\max_j|\phi_j(0)|,\qquad
 s=\max\{1,\sup_{j,|\operatorname{Im}z|\le a/2}|\phi'_j(z)|\},
 \qquad
 t_2=\max\{1,\sup_{j,|\operatorname{Im}z|\le a/2}|\phi''_j(z)|\}.
\]
Thus \(1+b,s,t_2,16/a\le\beta\), and the bound used for
activation values is \(|\phi_j(z)|\le b+s|z|\), not boundedness.
Write \(v=x/\sqrt d\).  Backward fields at a passive input are defined
by the same algebraic recursion as at a training input,
\[
 k^{(L)}(t,v)=w(t),\qquad
 \delta^{(j)}(t,v)=\phi'_j(z^{(j)}(t,v))\odot k^{(j)}(t,v),
 \qquad k^{(j)}=W^{(j+1)\top}\delta^{(j+1)}.
\]
This definition uses no passive label and does not add a training term.

\paragraph{1. Local coefficient estimates and working domains.}
\label{cp:src-coefficients}

The following scalars are bounds, not coordinates of a model:
\begin{align}
 H_1&=\max(1,b+20s),& H_j&=\max(1,b+10sH_{j-1}),\nonumber\\
 k_L&=H_L,& k_j&=10sk_{j+1},& \tau_j&=sk_j,\nonumber\\
 P_1&=3,&P_j&=H_{j-1}+10sP_{j-1}+1,&f_j&=sP_j,\nonumber\\
 A_*&=2sP_L+2s^2\sum_{j=2}^Lk_jP_{j-1},&
 D_*&=t_2\sum_{j=1}^LP_j^2,&
 H_*&=A_*+t_2\sum_{j=1}^LP_j^2k_j. \label{cp:src-basic}
\end{align}
Set \(f_* =\max_jf_j\), \(k_* =\max_jk_j\),
\(H_{\max}=\max_jH_j\), and define
\begin{align}
 E_*&=t_2\sum_{j=1}^L(10s)^{2(j-1)}k_j,&
 T_*&=2H_*^2+4A_*^2(e^2-1),\nonumber\\
 D_0&=(1+b+s)(1+T_*+E_*+s^2k_*^2),&
 D_1&=576e^3(1+b+s)D_*^3,\nonumber\\
 C_F&=s(8f_*^2+H_{\max}^2+1),&
 C_{\rm abs}&=8(D_0+1),\nonumber\\
 V_1^0&=\tau_1,&V_j^0&=\tau_jH_{j-1}^2+10sV_{j-1}^0,\nonumber\\
 G_L^0&=sH_L+14t_2V_L^0+s,&
 G_j^0&=10sG_{j+1}^0+s^3k_{j+1}^2H_j+14t_2V_j^0+s,\nonumber\\
 W_G&=128(1+H_{\max}+s\max_jV_j^0+\max_jG_j^0).
 \label{cp:src-absorption-coefficients}
\end{align}
Choose the exponential-budget constants in this order:
\begin{align}
 \eta_*&=\min\{1,(1024C_{\rm abs}W_G)^{-1}\},&
 \mathcal B&=1024e^2L,&
 \Lambda_*&=\log(e+\mathcal B)+\log(1/\eta_*).\nonumber
\end{align}
\begin{equation}
\begin{split}
 S_*={}&\min\left\{1,(4A_*)^{-1},
 \sqrt{\frac{\eta_*}{4000D_*}},
 \sqrt{\frac{\eta_*}{16eD_*\sqrt{2\mathcal B}}},\right.\\
 &\hspace{19mm}\left.\sqrt{\frac{\eta_*}{\Lambda_*}},
 \left(\frac{\eta_*^3}{D_1\mathcal B}\right)^{1/4},
 (8D_0C_F)^{-1/2}\right\}. \label{cp:src-allowance}
\end{split}
\end{equation}
These quantities depend only on the activation bounds and depth.
Further put
\begin{align}
 C_G&=32\max(1,H_{\max},\max_j\tau_j),&
 K_{\rm src}&=16C_{\rm abs}(1+C_G),\nonumber\\
 g_*&=\tau_1+\sum_{j=2}^L\tau_jH_{j-1},&
 r_j^*&=P_jg_*,&q_j^*&=f_jg_*,\nonumber\\
 j_j^*&=20(10s)^{j-1},&b_j^*&=sj_j^*,\nonumber\\
 e_1&=t_2r_1^*P_1,&
 e_j&=t_2r_j^*P_j+s(q_{j-1}^*+\tau_jH_{j-1}f_{j-1}+10e_{j-1}),\nonumber\\
 a_1^*&=t_2j_1^*P_1+2s,&
 a_j^*&=t_2j_j^*P_j+s(b_{j-1}^*+10a_{j-1}^*),\nonumber\\
 T_Q&=8f_*\max_j(f_jH_*+e_j),&T_J&=8f_*\max_ja_j^*.
 \label{cp:src-response-coefficients}
\end{align}
For a Gaussian multiplier \(G\), define
\begin{align}
 U_1(G)&=4sK_{\rm src},\nonumber\\
 U_j(G)&=2\{sK_{\rm src}(H_{j-1}^2+f_{j-1}^2+ST_Q
                  +S^2H_{j-1}q_{j-1}^*)+Gq_{j-1}^*+1\},\nonumber\\
 V_1(G)&=2(2G+2sK_{\rm src}S^2+1),\nonumber\\
 V_j(G)&=2\{Gb_{j-1}^*+sK_{\rm src}S^2
                    (T_J+H_{j-1}b_{j-1}^*)+1\}. \label{cp:src-UV}
\end{align}
Use \(U=\max_jU_j(16\sqrt{d+3})\),
\(V=\max_jV_j(16\sqrt{d+3})\) for the sphere, and
\(U_{\rm fin}=\max_jU_j(64)\) for a fixed finite list of inputs.
Finally let
\begin{equation}
 \mathcal K=H_L^2+S^2\left(\tau_1^2+
            \sum_{j=2}^L\tau_j^2H_{j-1}^2\right),\qquad
 D_W=\max\{\tau_1,\max_{j\ge2}\tau_jH_{j-1}\}.
 \label{cp:src-K}
\end{equation}

\begin{claim}[Local power bounds]
\label{cp:src-powers}
If \(Y/\lambda\le\beta^{-30L}\), then \(S\le\beta^{-26L}\le S_*\).
Moreover
\begin{gather}
 H_j\le3\beta^{2j},\quad P_j\le\beta^{3j-2},\quad
 k_j\le3\beta^{4L-2j},\quad\tau_j\le3\beta^{4L-2j+1},
 \quad K_{\rm src}\le\beta^{21L},\label{cp:src-powers-basic}\\
 U\le\beta^{26L}\sqrt{d+3},\qquad
 V\le\beta^{2L}\sqrt{d+3},\qquad
 U_{\rm fin}\le\beta^{26L},\qquad
 \mathcal K\le\beta^{14L}. \label{cp:src-powers-UV}
\end{gather}
\end{claim}
\begin{proof}
Unroll the affine recurrences using \(10s\le\beta^2\),
\(4j\le\beta^j\), \(L,L+1\le\beta^{L-1}\), and
\(\sum_{j=1}^L\beta^{cj}\le\beta^{cL}/(1-\beta^{-c})\).
They give \(H_j\le22\beta(10\beta)^{j-1}\le3\beta^{2j}\)
and \(P_j\le4j(10\beta)^{j-1}\le\beta^{3j-2}\).
Substitution gives the following bounds; both columns are included to
account for the numerical factors.
\begin{center}\small
\begin{tabular}{c|c|c}
quantity& intermediate bound&power bound\\\hline
\(A_*\)&\(7\beta^{5L-3}\)&\(\beta^{5L-2}\)\\
\(D_*\)&\(2\beta^{6L-3}\)&\(\beta^{6L-2}\)\\
\(H_*\)&\(4\beta^{8L-3}\)&\(\beta^{8L-2}\)\\
\(E_*\)&\(4\beta^{6L-3}\)&\(\beta^{6L-2}\)\\
\(T_*\)&\(3\beta^{16L-4}\)&\(\beta^{16L-3}\)\\
\(D_0\)&\(8\beta^{16L-3}\)&\(\beta^{16L-2}\)\\
\(D_1\)&\(9216e^3\beta^{18L-8}\)&\(\beta^{18L-2}\)\\
\(C_F\)&\(9\beta^{6L-1}\)&\(\beta^{6L}\)\\
\(C_{\rm abs}\)&\(65\beta^{16L-3}\)&\(\beta^{16L-1}\)
\end{tabular}\end{center}
For instance, the summands of \(A_*\)'s mixed part, \(H_*\)'s
curvature part, and \(E_*\) are at most
\(6\beta^{4L+j-3},3\beta^{4L+4j-3},3\beta^{4L+2j-3}\).
The inner sum in \(D_0\) is at most
\(1+3\beta^{16L-4}+4\beta^{6L-3}+9\beta^{8L-2}
\le4\beta^{16L-4}\), and \(9216e^3<\beta^6\).
Unrolling the two activity recurrences gives
\[
 V_j^0\le27j\beta^{4L+2j-3},\quad \max_jV_j^0\le\beta^{7L-2},
 \quad G_j^0\le66L\beta^{8L-2j-1},\quad
 W_G\le\beta^{9L+1}.
\]
In the reverse recurrence the propagated terminal, mixer and curvature
contributions are bounded respectively by
\(380L\beta^{8L-2j-2},27L\beta^{8L-2j-1},
379L\beta^{8L-2j-6}\); the constant contribution is at most
\(2\beta^{2L-2j-1}\).  Hence
\(\eta_*^{-1}\le\beta^{25L+4}\),
\(\mathcal B\le\beta^{L+3}\), and \(\Lambda_*\le\beta^{2L}\).
The six nonconstant entries of \eqref{cp:src-allowance} are at least
\(\beta^{-c}\), with exponents
\[
 5L-1,\quad(31L+6)/2,\quad(63L+12)/4,\quad
 (27L+4)/2,\quad(94L+13)/4,\quad(22L-1)/2.
\]
All are at most \(26L\).  Also \(16\beta^{-30L}\le\beta^{-26L}\).
The remaining recurrences give
\begin{gather*}
 g_*\le\beta^{5L-1},\qquad r_j^*\le\beta^{5L+3j-3},\qquad
 q_j^*\le\beta^{5L+3j-2},\\
 e_j\le2\beta^{5L+6j-4},\qquad a_j^*\le3\beta^{5j-2},
 \qquad T_Q\le\beta^{14L-3},\qquad T_J\le\beta^{8L-1}.
\end{gather*}
For \(e_j\), its three forcing exponents are at most
\(5L+6j-4,5L+3j-4,4L+3j-4\), with coefficient nine
in the last term; geometric propagation by \(\beta^2\) gives the
displayed factor two.  The same calculation gives the factor three
for \(a_j^*\).  Further
\(K_{\rm src}\le1552\beta^{20L-2}\le\beta^{20L+2}\).
Now \(ST_Q,S^2H_{j-1}q_{j-1}^*\le1\), and the bracket in
\eqref{cp:src-UV} is at most \(11\beta^{6L-8}\).
Thus for \(j\ge2\)
\[
 U_j(16\sqrt{d+3})\le22\beta^{26L-5}
       +32\sqrt{d+3}\beta^{8L-5}+2,
 \quad U_j(64)\le22\beta^{26L-5}+128\beta^{8L-5}+2.
\]
The first-layer bound is \(4\beta^{20L+3}\).  These give the
asserted bounds on \(U,U_{\rm fin}\).  In \(V_j\), the term
containing \(S^2\) is at most one, giving
\(V_j\le64\sqrt{d+3}\beta^{2L-2}+4\) for \(j\ge2\);
the first value is at most \(67\sqrt{d+3}\).
Finally substitute \eqref{cp:src-powers-basic} in
\eqref{cp:src-K}; \(S\le1\) and the same geometric sum give
\(\mathcal K\le\beta^{14L}\).  Every comparison above holds
already at \(L=2,\beta=10\), and improves as either increases.
\end{proof}

\paragraph{Working whole-sphere estimates.}
The estimates below in fact require only the conclusion of the fitting
lemma and \(S\le S_*\), both already implied by the headline assumptions.
For the remainder of this proof redefine the working radii by
\[
 T=32\ell/\lambda,\quad c_t=\frac a{64YSU},\quad
 c_q=\min\{1/8,a/(8V)\},\quad r_t=c_t/\sqrt\ell,
 \quad r_q=c_q/\sqrt\ell.
\]
For every fixed confidence, eventually with that confidence all forward
and backward fields extend holomorphically to a neighborhood of
\[
 ([-r_t,T+r_t]+i[-r_t,r_t])\times\mathcal Q_{r_q},\qquad
 \mathcal Q_r=\{v\in\mathbb C^d:v^\top v=1,
                           \|\operatorname{Im}v\|_2\le\sinh r\}.
\]
For \(d=1\), replace \(\mathcal Q_r\) by the two points \(\{-1,1\}\).
Holomorphy in the query is relative to the complex quadric.
Normalized first-layer and hidden operator norms are at most ten, and
\[
 \|h^{(j)}\|_2/\sqrt n\le H_j,\qquad
 \|k^{(j)}\|_2/\sqrt n\le Sk_j,\qquad
 \|\delta^{(j)}\|_2/\sqrt n\le S\tau_j.
\]
For training points,
\(\max|z_{a,i}^{(j)}|\le K_{\rm src}\sqrt\ell\) and
\(\max|k_{a,i}^{(j)}|\le SK_{\rm src}\sqrt\ell\).
In mobility coordinates
\(\Theta=(W^{(1)},\sqrt nW^{(2)},\ldots,\sqrt nW^{(L)},w)\),
put \(F_a=nf_n(v_a)\).  Then
\[
 \max_{a,j,i,v}|(D_\Theta z^{(j)}(v)\nabla_\Theta F_a)_i|
       \le SU\sqrt\ell,\qquad
 \max_{j,i}|\partial_\zeta z_i^{(j)}(u\cos\zeta+v\sin\zeta)|
       \le V\sqrt\ell
\]
for real orthonormal frames \((u,v)\) in the query tube.
All four source families
\[
 h^{(j)},\qquad W_0^{(j)}h^{(j-1)},\qquad\delta^{(j)},\qquad
 W_0^{(j+1)\top}\delta^{(j+1)},
\]
omitting boundary images, have
coordinate modulus at most \(M_0\sqrt n\), where
\(M_0=10\max_j(H_j,\tau_j)\).  The real fitting trajectory and its
endpoint are those of Lemma~\ref{cp:fit}.
The proof below also establishes the finite-query and expanding-domain
variants.  Its probability step is shared; the domain verifications are
given last.

\paragraph{2. Independent stops and deterministic differentiation.}
\label{cp:src-local}

Delete a fixed number \(u\) of neurons at one layer, preserving the
original normalization \(n\) in every rectangular matrix.  For each
deleted interior neuron, its initialized incoming row transpose \(y_i\)
and outgoing column \(x_i\) are independent \(N(0,I/n)\), independent
of the retained initialization and other omitted pairs.  At the first
layer the incoming row instead has covariance \(I_d\) and evaluates
only its own initial preactivation; at the top there is no outgoing
root.  Use a separate initial test for each cavity, and set its
reference coefficient paths identically zero if it fails.  All such
paths are defined before integrating an omitted root.

For a nested domain \(D_s\), stop each cavity at its own first exit
level.  Stop every training sample's budget
\begin{equation}
 \mathcal H_a=\frac1n\sum_{j,i}
 \exp\{\eta_*(Z_{a,i}^{(j)}+K_{a,i}^{(j)}/S)\},\quad
 Z_{a,i}^{(j)}=\sup_{D_s}|z_{a,i}^{(j)}|,\quad
 K_{a,i}^{(j)}=\sup_{D_s}|k_{a,i}^{(j)}|,
 \label{cp:src-budget}
\end{equation}
at \(\mathcal B\) for the full network and \(2\mathcal B\) for
cavities.  Also stop operators at ten, contour activity at
\(2\int\rho|dt|=S\), and imaginary preactivations at \(3a/8\)
and \(7a/16\), respectively.  Impose provisional coordinate and
response maxima twice the claimed full bounds and four times the
cavity bounds, as well as the passive size and row bounds used below.
These have strict margins on the successful prefix.

For rectangles use \(D_s=sD_1\) and coordinate clipping \(P_s\).
For the query tube also clip its imaginary great-circle coordinate,
retaining every real frame.  If a cavity stops at \(s_c\), extend a
coefficient by \(b_c(P_{s_c}z)\).  With
\(a_0(z)=\operatorname{clip}_{[0,T]}\operatorname{Re}z\),
\begin{equation}
 P_{s_c}a_0(z)=a_0(P_{s_c}z),\qquad
 b_c(P_{s_c}z)-b_c(P_{s_c}a_0(z))
   =b_c(P_{s_c}z)-b_c(a_0(P_{s_c}z)).
 \label{cp:src-clipped}
\end{equation}
Clipping is measurable in retained initialization, is Lipschitz, and
fixes the successful prefix.  Its real anchors traverse a monotone
interval and then freeze.  Different cavities are extended separately
before subtraction.  The expanding domain below has the same
properties with a uniformly bounded Lipschitz constant.  These
extensions need not be holomorphic outside their own domains.

Write \(\Theta_0=\Theta(0)\), \(\bar u=(\Theta-\Theta_0)/S\),
\(\bar F_a=F_a(\Theta_0+S\bar u,e_a)/S\),
\(\bar r_a=\bar F_a/n-y_a/S\), and put bars on backward
fields divided by \(S\).  Here \(e_a\) is an additive forward port.
The two actual ports of a deletion are
\(e_a=\sum_iW_{:,i}^{(j+1)}h_{a,i}^{(j)}\) and
\(\bar q_a=\sum_iW_{i,:}^{(j)\top}\bar\delta_{a,i}^{(j)}\).
The exact retained equation is
\begin{equation}
 \dot{\bar u}=-\frac2m\sum_a\bar r_a
   \left(\nabla_{\bar u}\bar F_a+
           D_{\bar u}h_a^{(j-1)\top}\bar q_a\right).
 \label{cp:src-retained}
\end{equation}
At the top add \(n^{-1}\sum_i\bar u_{w,i}h_{a,i}^{(L)}\)
to \(\bar r_a\), multiplying both terms.  The reverse port does
not enter the forward residual.  The gradient blocks of \(\bar F_a\)
are \(S\bar\delta_a^{(1)}v_a^\top\),
\(S\bar\delta_a^{(j)}h_a^{(j-1)\top}/\sqrt n\), and
\(h_a^{(L)}\).

The physical recurrences and linear activation growth give
\eqref{cp:src-basic}'s RMS bounds whenever the operator and activity
stops hold.  For augmented mobility/port directions \(U\), write
\(Z_a^{(j)}U=Dz_a^{(j)}[U]\).  The exact derivative recursion is
\begin{align}
 Z_a^{(1)}U&=U_{W^{(1)}}v_a+U_{e^{(1)}},\nonumber\\
 Z_a^{(j)}U&=U_{\sqrt nW^{(j)}}h_a^{(j-1)}/\sqrt n
       +W^{(j)}\operatorname{diag}(\phi'_{j-1})Z_a^{(j-1)}U
       +U_{e^{(j)}}. \label{cp:src-first-derivative}
\end{align}
It is bounded by \(P_j\); its activated version is bounded by \(f_j\).
The Hessian, with all activations evaluated at the reference state, is
\begin{align}
 D^2F_a[U,V]={}&U_w^\top\operatorname{diag}(\phi'_L)Z_a^{(L)}V
 +V_w^\top\operatorname{diag}(\phi'_L)Z_a^{(L)}U\nonumber\\
 &+\sum_j(Z_a^{(j)}U)^\top
       \operatorname{diag}(k_a^{(j)}\phi_j'')Z_a^{(j)}V\nonumber\\
 &+\sum_{j\ge2}\frac{\delta_a^{(j)\top}}{\sqrt n}
 \{U_{\sqrt nW^{(j)}}\operatorname{diag}(\phi'_{j-1})Z_a^{(j-1)}V
 +(U\leftrightarrow V)\}. \label{cp:src-Hessian}
\end{align}
For a rectangular matrix let
\(\|B\|_{p,n}=n^{-1/p}(\sum_i\sigma_i(B)^p)^{1/p}\).
The budget implies
\(\|k_a^{(j)}\|_p/n^{1/p}\le(S/\eta_*)p(2\mathcal B)^{1/p}\),
because \(x^p\le(p/e)^pe^x\).  Each noncurvature Hessian block
factors through at most \(n\) coordinates, and each curvature block
contains one carrier diagonal.  Matrix ideal inequalities in
\eqref{cp:src-Hessian} therefore give
\begin{align}
 \|D^2F_a\|_{p,n}&\le A_*+(SD_*/\eta_*)p(2\mathcal B)^{1/p},
 \qquad \|D^2F_a\|_{\rm HS}/\sqrt n\le H_*,\nonumber\\
 \|D^2F_a\|_{\rm op}&\le A_*+(SD_*/\eta_*)\log(2n\mathcal B).
 \label{cp:src-Schatten}
\end{align}
These bounds also cover \(D_\Theta\delta\), since
\(\partial_{e^{(j)}}F_a=\delta_a^{(j)}\).  Residual mixtures
preserve the bounds by \(\sum_a|r_a|/(m\rho)\le1\).

Suppose for now that products of the negative-Gram propagators on
ordered disjoint contour pieces have norm at most two, and
\(\rho/S\le\lambda/8\).  The variational generator is
\[
 -\frac2{mn}\sum_ag_ag_a^\top-\frac2m\sum_ar_aD^2F_a,
 \qquad g_a=\nabla_\Theta F_a.
\]
Its full propagator obeys
\begin{equation}
 \|J(t,s)\|\le2\exp\{SA_*+(S^2D_*/\eta_*)\log(2n\mathcal B)\}
 \le Cn^{\kappa},\qquad \kappa=1/4000.
 \label{cp:src-propagator}
\end{equation}
Here and only in qualitative local estimates, \(C\) and powers of
\(\ell\) may depend on the fixed problem and deletion count.
The coefficient of \(\log n\) is bounded by
\eqref{cp:src-allowance}; it has not been hidden in a width threshold.

\begin{claim}[Local maps, remainders, and passive observables]
\label{cp:src-local-lemma}
On the independent stops just specified, with contours of length
\(O(\ell)\), all first insertion maps have norm
\(n^\kappa\operatorname{poly}(\ell)\).  Their dependence on the
supremum norm of scalar controls is Lipschitz with that bound, and
their raw terminal/time/frame derivatives are bounded by
\(\sqrt n n^\kappa\operatorname{poly}(\ell)\).
On the uniform Gaussian event proved below, the retained full/cavity
differences in preactivations, scaled carriers, responses and angular
derivatives are at most \(n^{-1/30}\) per coordinate and at most
\(2n^{1/100}\) in Euclidean norm, for every fixed deletion count.
\end{claim}
\begin{proof}
For deterministic controls \(\alpha_{a,i},b_{a,i}\), use ports
\(e_a=\sum_ix_i\alpha_{a,i}\), \(\bar q_a=\sum_iy_ib_{a,i}\).
The three port derivatives of \eqref{cp:src-retained} at zero ports are
\[
 -\frac2{mn}\nabla_{\bar u}\bar F_a\bar\delta_a^{(j+1)\top},
 \qquad -\frac2m\bar r_a(D_{\bar u}\bar\delta_a^{(j+1)})^\top,
 \qquad -\frac2m\bar r_a(D_{\bar u}h_a^{(j-1)})^\top.
\]
The first term accounts for the changed residual.  Integrating them
against \eqref{cp:src-propagator}, using gradient norm
\(O(\sqrt n)\), response RMS, and contour length \(O(\ell)\),
gives an insertion-map bound \(\sqrt u n^\kappa O(\ell^2)\).
The differentiated forward/backward recursions are exactly
\begin{align*}
 z_{[1]}^{(j)}&=SV_{\sqrt nW^{(j)}}h^{(j-1)}/\sqrt n
                      +W^{(j)}h_{[1]}^{(j-1)}+e_{[1]}^{(j)},\\
 h_{[1]}^{(j)}&=\phi'_j\odot z_{[1]}^{(j)},\\
 \bar k_{[1]}^{(j)}&=W^{(j+1)\top}\bar\delta_{[1]}^{(j+1)}
       +SV_{\sqrt nW^{(j+1)}}^\top\bar\delta^{(j+1)}/\sqrt n
       +\bar q_{[1]}^{(j)},\\
 \bar\delta_{[1]}^{(j)}&=\phi'_j\odot\bar k_{[1]}^{(j)}
       +\phi_j''\odot\bar k^{(j)}\odot z_{[1]}^{(j)}.
\end{align*}
At the first layer replace the matrix product by \(SV_{W^{(1)}}v\);
at the top \(\bar k_{[1]}=V_w\).  A reference carrier costs
\(O(\ell)\) once in the inhomogeneous last term.  Homogeneous
propagation costs \(10s\).  Induction gives the asserted map bounds;
subtracting controls gives their Lipschitz bounds because the linear
variational equation is linear in these controls.  Differentiating
these displayed recursions introduces only the terms
\(\phi''z_t\),
\(\phi'''z_tz_{[1]}\bar k+\phi''z_{[1]}\bar k_t
+\phi''\bar k z_{[1],t}\), and ordinary product derivatives.
There is at most one raw reference coordinate derivative, bounded by
\(\sqrt n\operatorname{poly}(\ell)\).  This proves the derivative
claim without a width power depending on depth.  Cauchy's formula on
the strict strip bounds \(\phi'''\) by \(4\beta^2\).

Here are the nonlinear estimates needed to justify this linearization.
Decompose the augmented displacement into \((V,e)+(U,0)\), and let
\(N\) bound its linear Euclidean variations, \(d_0\) their coordinate
maxima, and \(u_0\) bound \(\|U\|\).  Put
\[
 R=d_0N+d_0u_0+u_0^2,\qquad P=(N+u_0)^2/\sqrt n.
\]
Assume \(N\ge1\), \(d_0,u_0\le1\), \(N+u_0\le\sqrt n\),
and \(u_0+R+P\le1\).  If \(E_g=g_1-g_0-Dg_0(V+U,e)\)
and \(g_i=\phi'(z_i)\) in the third line below, the exact identities are
\begin{align}
 E_{z^{(j)}}&=W_1^{(j)}E_{h^{(j-1)}}+
              S\Delta\bar u_{\sqrt nW^{(j)}}h_{[1]}^{(j-1)}/\sqrt n,
 \nonumber\\
 E_{\bar k^{(j)}}&=W_1^{(j+1)\top}E_{\bar\delta^{(j+1)}}+
 S\Delta\bar u_{\sqrt nW^{(j+1)}}^\top\bar\delta_{[1]}^{(j+1)}/\sqrt n,
 \nonumber\\
 E_{\bar\delta}&=g_1\odot E_{\bar k}
 +(g_1-g_0)\odot\bar k_{[1]}
 +\bar k_0\odot(g_1-g_0-\phi''(z_0)z_{[1]}).
 \label{cp:src-exact-remainders}
\end{align}
The first and top boundary remainders vanish.  Forward and backward
first derivatives are bounded by \(J_f=\beta^{10L}\) and
\(J_b=\beta^{20L}(1+M)\), where \(M\ge1\) is a common
provisional bound for the absolute preactivations and scaled carriers
of the reference and full paths, including the deleted-neuron values
that define the scalar controls. Such an envelope is
\(\operatorname{poly}(\ell)\) under the stated stops.
The localized product estimate is
\(\|z_{[1]}^{\odot2}\|_2\le3J_f^2R\).
Insert \(z_0+z_{[1]}\) into activation subtraction.  Taylor's formula
and \eqref{cp:src-exact-remainders}, summed by geometric propagation,
give
\begin{equation}
 \max(\|E_z\|_2,\|E_h\|_2)\le\beta^{30L}(R+P),\qquad
 \max(\|E_{\bar k}\|_2,\|E_{\bar\delta}\|_2)
       \le\beta^{60L}(1+M)(R+P).
 \label{cp:src-remainder-bounds}
\end{equation}
Indeed the two gate cross terms cost at most
\(\beta[3J_fJ_bR+\beta^{30L}(R+P)(d_0+J_bu_0)]\) and
\(M[\beta\|E_z\|+\beta^2\|z_{[1]}^{\odot2}\|/2]\);
the matrix cross term costs \(J_bP\).  No changed carrier is used
as a presumed reference bound.

The hidden gradient remainder is
\[
 \frac S{\sqrt n}
 [E_{\bar\delta}h_1^\top+
       \bar\delta_{[1]}(h_1-h_0)^\top+\bar\delta_0E_h^\top],
\]
bounded by \(\beta^{70L}(1+M)(R+P)\).
The Hessian identity along partial joining segments bounds the
residual difference by \(\beta^{13L}(N+u_0)/\sqrt n\) and its
linearization error by
\(\beta^{87L}(1+M)(N+u_0)^2/(2n)\).
For a product of residual and gradient, its exact remainder is
\(\bar r_0E_g+(\Delta\bar r-D\bar r_0\Delta)g_0
+\Delta\bar r(g_1-g_0)\).

For the reverse force one additionally needs a lower adjoint probe.
For an omitted root \(y_i\), include all derivatives of
\(y_i^\top h^{(j-1)}\) in the Gaussian coordinate event.
For \(q=\sum_iy_ib_{a,i}\), let \(Q=2u\beta M\),
\(P_q=u\beta Md_0\).  Subtracting its adjoint recursion gives
\[
 \max_j\|k_{q,1}^{(j)}-k_{q,0}^{(j)}\|_2
 \le\beta^{20L}(P_q+Q/\sqrt n)(N+u_0),
\]
since the changed gate multiplies the reference probe maximum, the
changed matrix costs \(Q(N+u_0)/\sqrt n\), and propagation costs
\(11s\).  Integrating the corresponding Hessian gives
\[
 \|(Dh_1-Dh_0)^\top q\|_2
 \le\beta^{50L}(P_q+Q/\sqrt n)[N+u_0+(N+u_0)^2].
\]
Combining these estimates in \eqref{cp:src-retained} yields the
vector-field remainder bound
\begin{equation}
 \beta^{240L}(1+u)(1+M)
 \{\bar\rho[R+d_0(N^2+Nu_0+u_0^2)]+(1+\bar\rho)P\},
 \qquad \bar\rho=\rho/S. \label{cp:src-force-remainder}
\end{equation}
This also holds on the complex joining segments: the sufficient gate
\(d_0+\beta^{10L}u_0+\beta^{30L}(R+P)<a/32\)
keeps the segments and auxiliary Taylor points inside \(15a/32\).
A first-exit argument proves the gate's use before Taylor expansion;
the factor \(\beta^{240L}\) pays for the Cauchy third-derivative
bound.  Scaling introduced no inverse power of \(S\): in the
effective network with readout \(\bar u_w\), the affine change of
variables multiplies hidden directions by \(S\le1\).

For every existing outgoing column, and for an incoming row at a
layer \(j\ge2\), integrating the omitted-parameter equations gives
\[
 \|\Delta W^{(j+1)}_{:,i}\|
       \le S^2(\max_r\tau_r)(b+sM)/\sqrt n,
 \qquad
 \|\Delta W^{(j)}_{i,:}\|
       \le S^2H_{\max}sM/\sqrt n.
\]
At the first layer the incoming-row estimate is instead
\(\|\Delta W^{(1)}_{i,:}\|\le sS^2M\).
There is no retained reverse port below that layer: this row's
motion is included in the deleted neuron's scalar controls, rather
than in the learned retained-port remainder. Thus every existing
learned port entering the retained equation still induces a force
at most
\(S^2\beta^{260L}(1+u)^2(1+M)^3(1+\bar\rho)/\sqrt n\).
The top residual offset is at most \(uM(b+sM)/n\); its product
with both retained terms has the same enclosing bound.

It remains to specify all passive quantities in the local claim.
For evaluated input \(q\) and training driver \(a\), let
\(C_q^{(j)}=D_\Theta h_q^{(j)}\),
\(\bar R_{qa}^{(j)}=S^{-1}D_\Theta z_q^{(j)}g_a\), and
\(\bar Q_{qa}^{(j)}=\phi'_j\odot\bar R_{qa}^{(j)}\).  Then
\begin{align}
 \bar R_{qa}^{(1)}&=\bar\delta_a^{(1)}v_a^\top q,&
 \bar R_{qa}^{(j)}&=\bar\delta_a^{(j)}c_{aq}^{(j-1)}
                         +W^{(j)}\bar Q_{qa}^{(j-1)},&
 c_{aq}^{(j-1)}&=h_a^{(j-1)\top}h_q^{(j-1)}/n.\label{cp:src-response-recursion}
\end{align}
Angular fields satisfy \(J^{(1)}=W^{(1)}q'\),
\(J^{(j)}=W^{(j)}I^{(j-1)}\), \(I=\phi'\odot J\).
Differentiate these product formulas: the four response terms are
\(\bar\delta_{[1]}c+\bar\delta c_{[1]}\),
\(SV_{\sqrt nW}\bar Q/\sqrt n\), and
\(W\bar Q_{[1]}\), together with the omitted port
\(\sum_ix_i\bar Q_{qa,i}\).  The activated variation adds
\(\phi''\bar R z_{[1]}\).  Angular variations have the same
matrix and gate terms, with omitted port \(\sum_ix_iI_i\).
Below a deletion there is also the direct reverse observable
\begin{equation}
 S^{-1}(Q_{qa,\mathrm{full}}-C_qg_a)=C_qC_a^\top\bar q_a.
 \label{cp:src-direct-reverse}
\end{equation}
Include the forward images of \(C_a^{0\top}y_i\) in the Gaussian
event; their coefficient maps are independent of that root.

Product subtraction gives
\[
 |c_{[1]}|\le\beta^{10L}(N+u_0)/\sqrt n,\qquad
 |E_c|\le\beta^{40L}[(R+P)/\sqrt n+(N+u_0)^2/n].
\]
The exact response remainder is
\(E_{\bar\delta}c_1+\bar\delta_{[1]}(c_1-c_0)
+\bar\delta_0E_c+W_1E_{\bar Q}+\Delta W\bar Q_{[1]}
+E_{\rm ports}\); its gate remainder is the third line of
\eqref{cp:src-exact-remainders} with \(\bar k\) replaced by
\(\bar R\).  The angular case uses that same identity with \(J\).
Reference response RMS cancels the \(1/\sqrt n\) in \(E_c\).
For \(\xi=C_a^{0\top}y_i\), forward directional subtraction gives
\(\|(C_q^1-C_q^0)\xi\|\le
\beta^{50L}(d_0+n^{-1/2})(N+u_0)\).
The other changed factor in \eqref{cp:src-direct-reverse} uses the
adjoint-probe estimate above.  Their sum is bounded by
\[
 \beta^{70L}(1+u)(1+M)(d_0+n^{-1/2})
                       [N+u_0+(N+u_0)^2].
\]
Missing feature products cost \(u(b+sM)^2/n\), and learned ports
cost \(n^{-1/2}\operatorname{poly}(\ell)\).  These enumerate all
additional remainders.  Their time and frame derivatives have the
same product structure already checked.  The probabilistic and
first-exit conclusion follows in the next paragraph.
\end{proof}

\paragraph{The uniform Gaussian event and closing the local comparison.}
For fixed controls, a pairing of an omitted root with a linear
variation has the form \(g^\top Rg\), where
\(R=(E_i^\top L+L^\top E_i)/2\).  Consequently
\(\|R\|\le\|L\|\), and only its same-root diagonal block has a
mean \(\operatorname{tr}R/n\).  Cross-root forms are centered.
On \(\|g\|\le2\sqrt{2u}\), interpolation of a centered form costs
at most \(10u\|\Delta R\|\), including its trace.
All stopped scalar controls have polynomial-logarithmic amplitudes
and speeds \(\sqrt n\operatorname{poly}(\ell)\); this follows by
differentiating the forward, backward and response recursions above.
A scalar control with amplitude \(M_c\), speed \(D_c\sqrt n\)
on a contour of length \(T_c\) has an \(n^{-1/8}\)-net of
logarithmic cardinality at most
\[
 (2+8T_cD_cn^{5/8})\log(1+8M_cn^{1/8}).
\]
Indeed sample it at spacing \(n^{-1/8}/(8D_c\sqrt n)\), round
on mesh \(n^{-1/8}/4\), and interpolate.  Apply this separately to
the finitely many real and imaginary controls.  They are deterministic
one-dimensional histories on a prescribed contour; no arbitrary
two-dimensional analytic continuation of controls is assumed.

After eventual absorption of the fixed logarithmic factors, all
linear and quadratic maps have norm at most \(n^{1/200}\).
A linear coordinate has variance at most \(n^{-99/100}\).
For a real symmetric form, diagonalization and
\[
 \mathbb E e^{t(G^\top RG-\operatorname{tr}R)}
 =\prod_i e^{-t\lambda_i}(1-2t\lambda_i)^{-1/2},\qquad
 -\log(1-v)-v\le\frac{v^2}{2(1-|v|)}
\]
give, for \(g\sim N(0,I/n)\),
\[
 \Pr\{|g^\top Rg-\operatorname{tr}R/n|>t\}
 \le2\exp\left[-c\min\left\{
       \frac{n^2t^2}{\|R\|_{\rm HS}^2},
       \frac{nt}{\|R\|_{\rm op}}\right\}\right]
\]
with an absolute \(c>0\), by applying exponential Markov with
\(|2t\lambda_i|\le1/2\) and optimizing the parameter.
Here \(\|R\|_{\rm HS}\le\sqrt{2un}\|R\|\).  At threshold
\(n^{-1/10}/8\), real/imaginary splitting gives
\(4e^{-n^{.78}}\) eventually.  Root-norm failure is at most
\(e^{-un}\).  The control interpolation error is
\(n^{-1/8+1/200}\operatorname{poly}(\ell)=o(n^{-1/10})\).
A terminal/time/frame mesh of spacing
\(n^{-2}/\operatorname{poly}(\ell)\) has polynomial size;
the sphere frame has \(2d\) ambient coordinates.  Together with
all coordinates, probes and at most \(uLn^u\) deletion sets, the
logarithmic union cost is eventually less than \(n^{.65}\).
Thus the uniform event fails with probability at most \(e^{-n^{.7}}\).
Only after this event is established are the actual adaptive controls
substituted.  Root-norm bounds also make the total linear vector
variation at most \(n^{1/100}\), distributing a fixed share among
the finitely many components.

Let \(U\) be the retained scaled displacement minus its linear
variation.  It satisfies Duhamel's formula with
\eqref{cp:src-force-remainder}, the learned ports, and the top offset.
Set \(N=n^{1/100}\), \(d_0=n^{-1/10}\), \(u_0=n^{-1/25}\).
On a first-exit prefix \(\|U\|\le u_0\),
\begin{equation}
 \|U\|\le n^\kappa\operatorname{poly}(\ell)
 [n^{-.08}+u_0^2+n^{-.48}+n^{-1/2}]=o(u_0).
 \label{cp:src-local-close}
\end{equation}
The activity integral bounds the residual-weighted terms and the
contour length bounds the others.  This strictly improves the stop.
The explicit remainders then give coordinate discrepancies
\(n^{-1/30}\) and vector discrepancies \(2n^{1/100}\), including
passive responses and angular fields.  The joining-strip gate also
holds eventually. For every deterministic frozen control history, the
retained linear variations are centered Gaussian images conditional on
the cavity. After substituting adaptive controls, we use only the
uniform bounds already proved; we assert neither Gaussianity nor
centering of the resulting variations. Pairing a frozen-control
variation again with an omitted root produces a quadratic form. Its
nonzero trace is retained before adaptive substitution, as treated next.
This completes Claim~\ref{cp:src-local-lemma}.

\paragraph{3. Trace absorption and Gaussian moments.}
\label{cp:src-moments}

\paragraph{Every noncentered contraction.}
Throughout this paragraph, conditional means and centered forms refer
to deterministic frozen controls, conditional on the cavity. The
resulting trace bounds hold uniformly over those controls and are
then evaluated at the actual adaptive controls.
For one interior omitted neuron work in unscaled mobility coordinates,
and set \(C_a=D_\Theta h_a^{(j-1)}\),
\(B_a=D_\Theta\delta_a^{(j+1)}\),
\(E_a=D_{e_a^{(j+1)}}\delta_a^{(j+1)}\), all at the zero-source
cavity.  With \(\alpha_b=h_{b,i}^{(j)}\),
\(b_b=\delta_{b,i}^{(j)}\), and \(d_b=\delta_b^{(j+1)}\),
the linear parameter variation is
\[
 V(t)=-\frac2m\sum_b\int_0^tJ(t,s)
 \{g_b(s)d_b(s)^\top x_i\alpha_b(s)/n
       +r_b(s)B_b(s)^\top x_i\alpha_b(s)
       +r_b(s)C_b(s)^\top y_i b_b(s)\}\,ds.
\]
The scalar expansions are
\[
 z_{a,i}=y_i^\top h_a^0+y_i^\top C_aV
                      +(\Delta W_{i,:})h_a+o(1),\qquad
 k_{a,i}=x_i^\top\delta_a^0+x_i^\top B_aV
                  +x_i^\top E_ax_i\alpha_a
                  +(\Delta W_{:,i})^\top\delta_a+o(S).
\]
The errors follow in scaled coordinates from
Claim~\ref{cp:src-local-lemma}.  The four nonzero conditional means are
\begin{align}
 &-\frac2m\sum_b\int r_bb_b
           \frac{\operatorname{tr}(C_aJ C_b^\top)}n\,ds,
 &&-\frac2m\sum_b\int r_b\alpha_b
           \frac{\operatorname{tr}(B_aJ B_b^\top)}n\,ds,\nonumber\\
 &\alpha_a\operatorname{tr}(E_a)/n,
 &&-\frac2{mn^2}\sum_b\int\alpha_b d_b^\top B_aJg_b\,ds.
 \label{cp:src-traces}
\end{align}
All other forms are centered same-root or independent-root forms
covered by the uniform Gaussian event.  The last displayed mean is
\(n^{-1+\kappa}\operatorname{poly}(\ell)=o(1)\), since
\(\|d_b\|,\|g_b\|=O(\sqrt n)\).

Expand \(J\) around its negative-Gram base.  A term with \(h\)
residual-Hessian insertions has ordered activity integral at most
\(S^h/h!\).  The product of its base propagators costs at most two.
Normalized Schatten H\"older, with exponent \(h+2\) on both
endpoint Hessian blocks and the \(h\) insertions, gives
\[
 \frac{2S^h}{h!}
 [A_*+(SD_*/\eta_*)(h+2)(2\mathcal B)^{1/(h+2)}]^{h+2}.
\]
The zero-insertion term is at most \(2H_*^2\).  Split the other
powers by \((x+y)^{h+2}\le2^{h+1}(x^{h+2}+y^{h+2})\).
The bounded part sums to \(4A_*^2(e^{2SA_*}-1)\); the carrier part is
\[
 \frac{8e^2D_*^2S^2\mathcal B}{\eta_*^2}
 \sum_{h\ge1}(2eD_*S^2/\eta_*)^h(h+2)^2
 \le\frac{576e^3D_*^3\mathcal B S^4}{\eta_*^3}.
\]
Indeed \(h!\ge(h/e)^h\), \((1+2/h)^h\le e^2\), and
\(\sum_{h\ge1}q^h(h+2)^2\le36q\) for \(q\le1/2\), by
differentiating the geometric series.  The allowance gives this last
condition.  For forward endpoints give one exponent two, the other
infinity, and the \(h\) Hessians exponent \(2h\).  The two series
are bounded by \(e^{2SA_*}-1\) and
\(\sqrt{2\mathcal B}\sum_{h\ge1}(4eD_*S^2/\eta_*)^h\),
each at most one by \eqref{cp:src-allowance}.  Including the zero term,
\begin{equation}
 |\operatorname{tr}(C_aJC_b^\top)|/n\le8f_*^2.
 \label{cp:src-forward-trace}
\end{equation}
Reciprocal exponents sum to one, giving normalization exactly
\(n^{-1}\), independent of parameter-space dimension.
The direct port Hessian is
\(E_a=\sum_{r>j}T_{r,j+1}^\top
\operatorname{diag}(k_a^{(r)}\phi_r'')T_{r,j+1}\), with
\(\|T_{r,j+1}\|\le(10s)^{r-j-1}\); hence carrier RMS gives
\(|\operatorname{tr}E_a|/n\le SE_*\).

Let \(G_{h,a,i}\), \(G_{\delta,a,i}\) denote suprema of the
corresponding cavity Gaussian pairings.  Set
\(Z_i=(m^{-1}\sum_aZ_{a,i}^2)^{1/2}\) and
\(U_i=(m^{-1}\sum_a(K_{a,i}/S)^2)^{1/2}\).
At every driving time Cauchy--Schwarz gives
\(m^{-1}\sum_b|r_bh_{b,i}|\le\rho(b+sZ_i)\) and
\(m^{-1}\sum_b|r_b\delta_{b,i}|\le\rho sSU_i\).
Integrating the rank-one updates and pairing retained factors in RMS,
the learned outgoing column costs \(s^2k_*^2S^3\) times its driving
amplitude, and the learned incoming row costs
\(sS^2U_iH_{j-1}^2\); at the first layer the latter is \(sS^2U_i\).
Thus \eqref{cp:src-traces} gives
\begin{align}
 K_{a,i}/S&\le G_{\delta,a,i}/S+
 (D_0+D_1\mathcal BS^4/\eta_*^3)(1+Z_{a,i}+Z_i)+o(1),\nonumber\\
 Z_{a,i}&\le G_{h,a,i}+C_FS^2U_i+o(1).
 \label{cp:src-scalar-feedback}
\end{align}
The direct external trace keeps the evaluated sample's amplitude
\(Z_{a,i}\); only integrated terms use the driving-sample RMS.
At the first layer \(G_{h,a,i}=|W_{0,i}^{(1)}v_a|\); at the top
\(\sup|w_i|/S\le b+sZ_i\) and there is no outgoing root.
These endpoint formulas obey the same enclosing inequalities.

Put \(\overline G_h=(m^{-1}\sum_aG_{h,a,i}^2)^{1/2}\) and
\(\overline G_\delta=(m^{-1}\sum_a(G_{\delta,a,i}/S)^2)^{1/2}\).
The backward coefficient is at most \(2D_0\) and
\(D=C_FS^2\le1/(8D_0)\).  Taking sample RMS gives
\[
 U_i\le\overline G_\delta+2D_0(1+2Z_i)+o(1),\qquad
 Z_i\le\overline G_h+DU_i+o(1).
\]
Since \(4D_0D\le1/2\), solve these two inequalities and substitute
back sample by sample to obtain
\begin{equation}
 Z_{a,i}+K_{a,i}/S\le C_{\rm abs}
 (1+G_{h,a,i}+G_{\delta,a,i}/S+
                     \overline G_h+\overline G_\delta)+o(1).
 \label{cp:src-scalar-absorption}
\end{equation}
This coefficient is independent of sample and deletion counts.

\paragraph{Real reference moments.}
Parameterize a real stopped cavity by its own normalized activity
\(u=2\int_0^t\rho/S\in[0,1]\), freezing after its endpoint.
For \(v=|u-u'|\), integrating the updates gives
\begin{gather*}
 \|\Delta w\|_2/\sqrt n\le H_LSv,\qquad
 \|\Delta W^{(1)}\|_F/\sqrt n\le\tau_1S^2v,\\
 \|\Delta W^{(j)}\|_F\le\tau_jH_{j-1}S^2v,\qquad
 \|\Delta z_a^{(j)}\|_2/\sqrt n\le V_j^0S^2v.
\end{gather*}
The feature difference has the last bound multiplied by \(s\).
Split a reference carrier at \(SR\).  Its high part satisfies
\(\|k\mathbf1_{|k|>SR}\|_2/\sqrt n
\le(4S/\eta_*)\sqrt{2\mathcal B}e^{-\eta_*R/4}\), by the
budget and \(x^2e^{-x/2}\le16\).
Choose \(R=(4/\eta_*)\log[8\sqrt{2\mathcal B}/(\eta_*v)]\)
when \(v>0\).  The high changed-gate part is at most \(sSv\);
the low part is at most \(t_2V_j^0S^3Rv\).  Since
\(R\le\eta_*^{-1}(10\Lambda_*+4\log(1/v))\),
\(S^2\Lambda_*/\eta_*\le1\), and
\(v\log(1/v)\le2\sqrt v/e\), the low part is bounded by
\(14t_2V_j^0S\sqrt v\).  At the top add \(sH_LSv\);
below it add the changed-matrix term
\(s^3k_{j+1}^2H_jS^3v\) and propagated difference.  This proves
\begin{equation}
 \|\delta_a^{(j)}(u)-\delta_a^{(j)}(u')\|_2/(S\sqrt n)
                 \le G_j^0\sqrt{|u-u'|}.
 \label{cp:src-real-modulus}
\end{equation}
The top carrier was included; no coordinate readout bound was used.

Conditioned on the cavity, a Gaussian reference has bounded initial
variance and this square-root activity modulus.  Grids of spacing
\(4^{-q}\) have at most \(2\cdot4^q\) increments of standard
deviation at most \(2G2^{-q}\).  Gaussian tails at thresholds
proportional to \(2^{-q}\sqrt{z^2+4(q+1)}\), summed over levels
and the initial value, give
\[
 \Pr\{X>W_G(1+z)\mid\text{cavity}\}\le4e^{-z^2},\qquad
 \mathbb E[e^{X^2/(16W_G^2)}\mid\text{cavity}]<2.
\]
The factor 128 in \(W_G\) bounds the geometric threshold sums.
For the moment, \((1+z)^2\le2(1+z^2)\) and tail integration give
\(\mathbb E e^{Z^2/8}\le1+4/7\), with
\(e^{1/8}(1+4/7)<2\).  Jensen gives the same bound for sample RMS:
\(e^{\overline X^2/(16W_G^2)}\le
m^{-1}\sum_ae^{X_a^2/(16W_G^2)}\), without sample independence.
Completing the square and H\"older for the four terms in
\eqref{cp:src-scalar-absorption} give the one-neuron moment base
\begin{equation}
 2\exp\{qC_{\rm abs}+64q^2C_{\rm abs}^2W_G^2\}<4,
                    \qquad 0\le q\le8\eta_*.
 \label{cp:src-main-moment}
\end{equation}
The chosen \(\eta_*\) verifies the last inequality.

\begin{claim}[Local Gaussian rectangle bound]
\label{cp:src-Gaussian}
Let a deterministic complex coefficient map on a rectangle of sides at
most \(H\) satisfy \(\sup_z\|b_z\|\le D\) and
\(\|b_z-b_{z'}\|\le K\|z-z'\|\).  For a standard real Gaussian
vector \(G\), put \(X=\sup_z|G^\top b_z|\).  Then
\[
 M=256D\sqrt{\log(e+HK/D)},\qquad
 \log\mathbb E e^{qX}\le qM+2q^2D^2\quad(q\ge0).
\]
For sample RMS with the same bounds, replace \(M\) by \(M+2D\).
If \(D=0\), the process is zero.  For covariance \(I/n\), divide
coefficient norms by \(\sqrt n\).
\end{claim}
\begin{proof}
Use dyadic coefficient meshes with spacing proportional to
\(D2^{-j}/K\), having at most \(4(1+HK/D)^2 4^j\) points.
Connect each point to the preceding mesh; increment standard deviations
are at most \(3D2^{-j}\).  Gaussian exponential moments give
\(\mathbb E\max_{r\le N}|Z_r|\le\sigma\sqrt{2\log(2N)}\)
when standard deviations are at most \(\sigma\).  Summing uses
\(\sum_{j\ge1}2^{-j}=1\), \(\sum_{j\ge1}2^{-j}\sqrt j\le2\).
A factor 128 covers the real process, and real/imaginary splitting
gives 256.

An \(a\)-Lipschitz Gaussian function satisfies
\(\log\mathbb E e^{t(F-\mathbb EF)}\le a^2t^2/2\) and
\(\operatorname{Var}F\le a^2\).  Here is a derivation.
For \(P_tg(x)=\mathbb E g(e^{-t}x+\sqrt{1-e^{-2t}}G)\),
differentiating invariant entropy and Gaussian integration by parts give
\[
 \operatorname{Ent}_\gamma(g)
 =\int_0^\infty\mathbb E\frac{|\nabla P_tg|^2}{P_tg}\,dt
 \le\tfrac12\mathbb E\frac{|\nabla g|^2}{g}.
\]
Use \(\nabla P_tg=e^{-t}P_t\nabla g\) and weighted
Cauchy--Schwarz for the inequality.  Apply it to \(g=e^{uF}\) and
integrate the resulting derivative inequality for
\(u^{-1}\log\mathbb E e^{uF}\).  Truncation and smoothing extend
the calculation to Lipschitz functions; their linear growth supplies
exponential integrability.  Differentiating at zero gives the variance.
Here \(X\) is at most \(2D\)-Lipschitz, proving the estimate.
The sample RMS is also \(2D\)-Lipschitz; its mean is at most
\(M+2D\) by the variance bound and Cauchy--Schwarz.  Independence
among sample processes is unnecessary.
\end{proof}

\paragraph{Initial cavities and maximum improvements.}
The initialization event of Lemma~\ref{cp:fit} applies jointly to
fixed-size rectangular cavities.  To see why weak covariance estimates
need not be union-bounded over deletion sets, zero-embed their features.
On the full initial operator event, conditional fresh-row Gaussian
tails give initial training maxima \(Z_0=C\sqrt\ell\) with
probability tending to one.  Deleting \(u\) coordinates changes each
later initialized feature vector by at most
\[
 D_{\rm init}=(8s)^L\sqrt u(b+sZ_0)=O_u(\sqrt\ell).
\]
Normalized feature and singular-value changes therefore vanish
uniformly over deletion sets, transferring the strict fitting margins.
For coordinate localization above a deletion, project the complete
lower feature difference onto this deterministic ball before pairing
with the independent current row.  At threshold \(n^{-1/10}\),
the tail \(2\exp[-n^{4/5}/(2D_{\rm init}^2)]\) pays for the
\(O_u(n^{u+1})\) tests.  A cavity's covariance retains normalization
\(n\); it is not identified with the full covariance.
Initial budgets are below \(\mathcal B/2\) with probability tending
to one: their conditional means are below
\(2e^{\eta_*^2H_{\max}^2/2}<3\), their second exponential
moments are bounded, and conditional Chebyshev applies layer by layer.
The coordinate-small deletion comparison transfers this event to all
cavities.  Their own initial tests remain part of each globally
extended reference's definition.

Under the provisional stops, Gaussian coefficient RMS is bounded by
\(H_{\max}\) or \(S\max_j\tau_j\).  A mesh of spacing \(n^{-2}\)
has polynomial size, with vanishing off-grid error from the normalized
coefficient derivative bounds.  Complex Gaussian splitting has tail
\(4e^{-u^2/(4\sigma^2)}\).  The multiplier \(C_G\), followed by
\eqref{cp:src-scalar-absorption}, improves the full training stops to
\begin{equation}
 \max|z_{a,i}^{(j)}|\le K_{\rm src}\sqrt\ell,
 \qquad\max|k_{a,i}^{(j)}|\le SK_{\rm src}\sqrt\ell.
 \label{cp:src-improved-maxima}
\end{equation}
There are at most \(n^7\) tests eventually on the original rectangle
and \(n^{10}\) on the expanding finite-panel rectangle; the factor
32 leaves strict margin.  This step uses Gaussian tails on a mesh,
not already removed budgets or Gaussian supremum moments.

For response and angular stops, differentiate \(Q_{qa}=C_qg_a\):
\(DQ_{qa}=C_qDg_a+(DC_q)[\cdot]g_a\).
The first term has normalized Hilbert--Schmidt norm at most \(f_jH_*\).
Curvature, changed matrix, changed feature and propagation in the
second give respectively
\(t_2r_j^*P_j,sq_{j-1}^*,s\tau_jH_{j-1}f_{j-1},10se_{j-1}\).
These are the recurrence for \(e_j\).  Angular fields give the
recurrence for \(a_j^*\), including the first-layer map
\(U_{W^{(1)}}q'\) of normalized norm at most two.
Apply \eqref{cp:src-forward-trace} with these endpoints.  The four
noncentered response terms are exactly the direct feature pairing,
the direct reverse term \eqref{cp:src-direct-reverse}, the integrated
mixed endpoint, and the learned incoming row.  After division by \(S\),
their coefficients are
\[
 sK_{\rm src}H_{j-1}^2,\quad sK_{\rm src}f_{j-1}^2,\quad
 sK_{\rm src}ST_Q,\quad sK_{\rm src}S^2H_{j-1}q_{j-1}^*.
\]
The centered incoming pairing contributes \(Gq_{j-1}^*\).
Angular terms are \(Gb_{j-1}^*\) and
\(sK_{\rm src}S^2(T_J+H_{j-1}b_{j-1}^*)\).
Thus the factor two in \eqref{cp:src-UV} strictly improves these stops.
For sphere frames, at most \(n^{4d+10}\) mesh points require
\(G=16\sqrt{d+3}\); polar normalization of nearby frames gives
bounded frame derivatives for interpolation.  A fixed finite input
list has no frame mesh and at most \(n^{10}\) tests eventually;
\(G=64\) gives failure \(4n^{10}e^{-1024\ell}\).
The number of passive inputs affects only the qualifying width.

Use the provisional passive-preactivation cap
\(2(GH_{\max}+1+U_0S^2+a)\ell^2\), where \(U_0\) is the
applicable response bound.  Its initial maximum is
\((GH_{\max}+1)\sqrt\ell\), by the same fresh-row argument;
real activity adds at most \(U_0S^2\sqrt\ell\).
The domain checks below control the complex displacement and improve
this stop without a passive exponential budget.

\begin{claim}[Local time derivative under provisional stops]
\label{cp:src-time-derivative}
Let \(U_0=U\) or \(U_{\rm fin}\), as appropriate, and put
\(N_* =\max_j\{r_j^*,4U_0\}\),
\(B_n=1+(S^2/\eta_*)\log(e+\ell)\).  Define
\[
 J_L=2sH_L+4et_2N_*,\qquad
 J_j=10sJ_{j+1}+2s^3k_{j+1}^2H_j+4et_2N_*,\qquad J_* =\max_jJ_j.
\]
For \(\ell\ge\max(e^2,2\mathcal B)\), on each independently
stopped cavity,
\[
 \|\dot h_a^{(j)}\|_2/\sqrt n\le2\rho Sq_j^*,\qquad
 \|\dot\delta_a^{(j)}\|_2/\sqrt n\le J_*B_n\rho.
\]
These bounds use the provisional response stops, not their improvement.
\end{claim}
\begin{proof}
The exact response identity gives
\(\|\dot z\|_2/\sqrt n\le2\rho Sr_j^*\),
\(\|\dot z\|_\infty\le8\rho SU_0\sqrt\ell\).
Interpolate to exponent \(2v/(v-2)\) and use the carrier
\(v\)-norm supplied by \eqref{cp:src-budget}.  For \(v\ge4\),
\[
 \|k\odot\dot z\|_2/\sqrt n
 \le(2\rho S^2N_*/\eta_*)v(2\mathcal B\ell)^{1/v}.
\]
Choose \(v=\max(4,\log(2\mathcal B\ell))\).  The exponential
factor is at most \(e\), and \(v\le2\log(e+\ell)\).
Backward differentiation now gives the readout term \(2sH_L\),
propagated derivative, changed-matrix term
\(2s^3k_{j+1}^2H_j\), and changed-gate term \(4et_2N_*\).
They are the displayed recursion.  Forward differentiation gives the
first bound.  This growing real exponent is a deterministic budget
inequality and does not use growing deletion sets.
\end{proof}

\paragraph{Removing budgets on complete independently stopped paths.}
The domain checks below will supply a bounding rectangle of side
\(O(\ell)\), uniformly Lipschitz retractions commuting with real
anchors, and normalized complex-minus-real coefficient radius
\(D_n=O(B_n/\sqrt\ell)\) with global modulus \(O(B_n)\).
Claim~\ref{cp:src-Gaussian} then gives mean
\(O((\log(e+\ell))^{3/2}/\sqrt\ell)\) and vanishing tail radius.
Every fixed exponential moment of these corrections tends to one.
Fixed squared-exponential moments do also, by integrating the
sub-Gaussian tail; sample RMS uses Jensen.  Notice that the modulus
here is polylogarithmic, not an arbitrary polynomial in \(n\).

For singleton/common-cavity differences, first extend each path by
its own retraction, then project their difference onto the deterministic
ball of radius \(4n^{1/100}\).  Both paths omit the paired root.
Projection preserves this independence and equals the actual difference
on the successful prefix.  Its normalized Gaussian radius is
\(4n^{-49/100}\), and its modulus is at most a fixed power of
\(n\), by the already proved derivative bounds.  The Gaussian
rectangle claim again gives moments tending to one at every fixed
order.  The two coefficient paths need not be independent of each other.

Fix a positive integer moment order \(u\), a sample and a layer.
Distinct roots in an empirical \(u\)th moment are independent given
their common cavity.  Apply Cauchy--Schwarz to separate products of
main real references from products of corrections, and then H\"older
over the finitely many roots and four reference families for the
corrections.  Main moments use exponent at most \(8\eta_*\) in
\eqref{cp:src-main-moment}; corrections use fixed orders proportional
to \(u\) and tend to one.  The limiting per-root base is therefore
below sixteen.  Drop full-event indicators only after dominating by
these nonnegative, globally defined cavity bounds.  A tuple with
\(k<u\) distinct roots has normalized multiplicity
\(O_u(n^{k-u})\).  On \eqref{cp:src-improved-maxima}, repeated
weights cost at most \(e^{2u\eta_*K_{\rm src}\sqrt\ell}=n^{o(1)}\),
so collision contributions vanish.  Local exceptional events are
superpolynomially small, and the stopped budgets are bounded.

Cavities cannot stop first on the successful full prefix: their
coordinate differences are \(n^{-1/30}\), maxima and poles have
strict transferred margins, and
\[
 \mathcal H_a^{\rm cavity}\le
 e^{2\eta_*n^{-1/30}}\mathcal H_a+O_u(1/n)<2\mathcal B.
\]
Operator and activity margins are established in the domain checks.
If a full budget is hit, some sample/layer mean is at least
\(\mathcal B/L\).  Markov's inequality and the fixed sample/layer
union give
\begin{equation}
 \limsup_{n\to\infty}\Pr\{\text{a full budget is hit}\}
              \le mL(16L/\mathcal B)^u. \label{cp:src-budget-removal}
\end{equation}
Take the width limit at each fixed \(u\), then the infimum over
positive integers.  Since \(16L/\mathcal B<1\), the result is zero.
The label allowance precedes the moment order and does not shrink
with confidence or sample count.  This is the shared probability step.

\paragraph{4. Verification of the working domains.}

\begin{proof}[Verification of the working whole-sphere estimates]
Lemma~\ref{cp:fit} independently supplies the real trajectory, strict
operator caps below nine, and real activity at most \(S/2\).
The same holds for the independently tested cavities.  Impose the
eventual deterministic gates
\begin{equation}
 n^{-1}\le\min(Y,S,1),\qquad \ell\ge\max(e^2,2\mathcal B),
 \qquad \sqrt\ell\ge c_t
       \max\{8,\lambda,4\mathcal K/\log2,32YSD_W\}.
 \label{cp:src-rectangle-gates}
\end{equation}
Follow the real solution to its nearest anchor in \([0,T]\), then
at most two short contour pieces of total length \(2r_t\).
The algebraic residual Gram and the negative-Gram variational
generator, divided by their factor two, have norm at most
\(\mathcal K\), by the explicit gradient blocks.  No complex
positivity is used.  Residual growth and products of the base
propagators cost at most \(e^{4\mathcal Kr_t}\le2\).
Additional activity is at most \(8Yr_t\le S/2\);
normalized hidden increments are at most \(8YSD_Wr_t\le1/4\).
Thus operators have strict margins inside ten and \(\rho/S\le\lambda/8\).
The total contour length is \(O(\ell)\).

These are the premises of the local insertion and trace argument,
which improves maxima and response/angular caps.  The exact velocity
\(\partial_tz^{(j)}=-(2/m)\sum_ar_aR_a^{(j)}\) now bounds the
complex-time preactivation displacement by \(8c_tYSU=a/8\).
Every query in the intrinsic tube has the form
\(u\cosh h+iv\sinh h\) for a real orthonormal frame and
\(|h|\le r_q\).  Its additional displacement is at most
\(c_qV\le a/8\).  Thus the full pole stop improves strictly;
the local coordinate comparison transfers this margin to cavities.
The passive size stop also improves as already shown.

Under the provisional response stops, Claim~\ref{cp:src-time-derivative}
bounds normalized complex-minus-real reference radii by
\(C r_tB_n=O(B_n/\sqrt\ell)\), and their global time moduli
by \(CB_n\).  This holds on complete independently stopped paths:
the clipped arguments commute with the real-anchor map as in
\eqref{cp:src-clipped}.  Their parameter rectangle has side
\(O(\ell)\).  The budget-removal argument therefore applies,
with probability of failure tending to zero.  The logical order is
provisional deterministic bounds, local insertion, maximum/response
improvement, strict pole/cavity transfer, and then complete-path
Gaussian moments and budget removal.  None of these steps conditions
an omitted Gaussian root on full-network survival.

Before each stop, finite-width bounded parameters and strict separation
from the activation singularities give local holomorphic ODE
continuation.  Compactness and uniqueness glue the local continuations
on the nested rectangles.  At a first exit the strict derivative strip
allows one-sided estimates; the strict improvements continue the
solution across the closed boundary.  Query holomorphy follows by
forward composition on the quadric.  Passive backward fields are then
holomorphic by their finite downward recursion, with RMS bounds
\(Sk_j,S\tau_j\), using the same readout and operator bounds;
they require no new probabilistic budget.  Initialized operator norm
eight proves the four-family coordinate bound.  All failures tend to
zero, hence each prescribed fixed confidence is attained eventually.
\end{proof}

\begin{claim}[Working finite-query time domain and all-time carriers]
\label{cp:src-finite-working}
Under \(Y/\lambda\le\beta^{-30L}\), for any fixed finite list of
real unit inputs, eventually with every fixed confidence its predictions
are holomorphic on a neighborhood of
\[
 [-r,T+r]+i[-r,r],\qquad
 T=32\ell/\lambda,\qquad r=1/(\lambda\sqrt\ell),
\]
and \(|f_n(t,v)|\le SH_L^2\) there.  Two independent copies thus
satisfy \(|f_n-\widetilde f_n|\le32\beta^{6L}Y/\lambda\)
on their common event.  In particular the Bernstein ellipse of
parameter two on \([0,r]\) is contained in this rectangle.
On the real all-time trajectory, also
\begin{equation}
 \sup_{t\in[0,\infty]}\max_{a,j,i}|k_{a,i}^{(j)}(t)|
               \le2SK_{\rm src}\sqrt\ell.
 \label{cp:src-carrier-working}
\end{equation}
The fitted endpoint is included.
\end{claim}
\begin{proof}
Use the same proof without angular fields or the frame mesh, replacing
the Gaussian multiplier by 64.  The finite-list response bound is
\(SU_{\rm fin}\sqrt\ell\).  The time-radius coefficient may be any
\(c\le a/(64YSU_{\rm fin})=a/(4\lambda S^2U_{\rm fin})\).
By Claim~\ref{cp:src-powers},
\[
 \frac a{4S^2U_{\rm fin}}
 \ge\frac{\beta^{34L-1}}{64}>1,
\]
so \(c=1/\lambda\) is permitted.  All other rectangle gates are
eventual.  Cauchy--Schwarz gives \(|f_n|\le SH_L^2\); use
\(H_L\le\beta^{3L}\) for the difference bound.  The parameter-two
ellipse on \([0,r]\) has real range \([-r/8,9r/8]\) and imaginary
height \(3r/8\), hence lies inside the stated rectangle eventually.

It remains to extend the carrier maximum beyond \(T\); no complex
late-time extension is required.  Let \(\|\Delta\theta\|\) denote
the physical parameter norm with blocks
\((W^{(1)}/\sqrt n,W^{(2)},\ldots,W^{(L)},w/\sqrt n)\).
Forward subtraction in the real fitting tube gives feature RMS
changes bounded by \(P_j\|\Delta\theta\|\).  Backward subtraction
then gives
\[
 \max_{a,j}\|k_a^{(j)}(t)-k_a^{(j)}(T)\|_\infty
 \le C_k n\|\theta(t)-\theta(T)\|,
\]
where \(C_L^k=1\),
\(C_j^k=S\tau_{j+1}+10sC_{j+1}^k+
10t_2SP_{j+1}k_{j+1}\), and \(C_k=\max_jC_j^k\le\beta^{20L}\).
The terms are the changed mixer, propagated upper difference, and
changed activation gate; RMS-to-coordinate conversions account for
the two factors \(\sqrt n\).  The real segment has bounded
operators and real activations, so this subtraction needs no complex
strip hypothesis.  Lemma~\ref{cp:fit} gives
\(\|\theta(t)-\theta(T)\|\le2Y e^{-16\ell}/\sqrt\lambda\).
Consequently the last display is eventually at most
\(SK_{\rm src}\sqrt\ell\), including at \(t=\infty\).
Add \eqref{cp:src-improved-maxima} at \(T\).
\end{proof}

\begin{claim}[Working residual-adapted finite-panel domain]
\label{cp:src-panel-working}
Fix the training inputs and any finite list of passive inputs before
initialization.  Use the same allowance as the working whole-sphere estimates,
and let
\begin{gather*}
 \nu=\lambda/4,\quad T=32\ell/\lambda,\quad
 r_n=\frac a{64YSU_{\rm fin}\sqrt\ell},\\
 h(t)=\frac1{2\mathcal K}\log(1+2\mathcal K r_ne^{\nu t}),
 \quad d_h=\frac\nu{2\mathcal K},\quad T_+=T+h(T)+r_n.
\end{gather*}
Eventually with every fixed confidence, the required forward and
backward fields are holomorphic on a neighborhood of
\[
 D_1=\{x+iy:-h(0)\le x\le T_+,\quad |y|\le h(\max(x,0))\},
\]
with the same operator, RMS, training-maximum, and four-family
coordinate bounds as in the working whole-sphere estimates; the response
bound is \(SU_{\rm fin}\sqrt\ell\).  In particular it suffices
to retain forward families on all declared inputs and backward
families only on training inputs.
For \(0\le t\le T\), the disk of radius
\(\varrho(t)=h(t)/(2(1+d_h))\) about \(t\) lies in \(D_1\).
The adaptive subdivision
\(t_0=0\), \(t_{b+1}=\min(T,t_b+\varrho(t_b)/2)\)
has at most
\begin{equation}
 J\le1+40960\frac{U_{\rm fin}}a
                  \left(\frac Y\lambda\right)^2\sqrt\ell
       +80\frac{\mathcal K}{\lambda}
                  \log\left(1+\frac{4\ell}{\log2}\right)
 \label{cp:src-panel-count}
\end{equation}
panels.  Here \(U_{\rm fin}/a\le\beta^{27L}\) and
\(\mathcal K\le\beta^{14L}\), independently of the input dimension
and the size of the fixed passive list.
\end{claim}
\begin{proof}
Population RMS and the Gram gap give
\(\lambda\le H_L^2\le\mathcal K\).  Thus
\(0\le h'\le d_h\le1/8\), \(h(0)\le r_n\), and
\(h(t)\le r_n+d_ht\).  Define the nested domains by
\[
 D_s=\{x+iy:-s h(0)\le x\le sT_+,\quad
                      |y|\le s h(\max(x,0))\},\quad0\le s\le1.
\]
For \(X_s(x)=\operatorname{clip}_{[-s h(0),sT_+]}x\), their
retractions are
\[
 P_s(x+iy)=X_s(x)+i\operatorname{clip}_{[-s h(X_s(x)_+),
                                           s h(X_s(x)_+)]}y.
\]
With \(a_0(z)=\operatorname{clip}_{[0,T_+]}\operatorname{Re}z\),
\(P_sa_0=a_0P_s\), and \(\operatorname{Lip}(P_s)\le1+d_h\).
Indeed \(|\Delta X_s|\le|\Delta x|\) and the imaginary
increment is at most \(|\Delta y|+d_h|\Delta x|\); the shear
has norm at most \(1+d_h\).  Integrate derivatives along the image
under \(P_s\) of a straight segment to get global coefficient
Lipschitz bounds.  Convexity of \(D_s\) is unnecessary.  The real
anchors remain monotone and freeze, as required by the moment proof.

We verify the base-propagator premise before Gaussian insertion.
Under each cavity's provisional response stop, use
Claim~\ref{cp:src-time-derivative} with \(U_0=U_{\rm fin}\).
For
\(G=(mn)^{-1/2}[\nabla_\Theta F_1,\ldots,\nabla_\Theta F_m]\),
the explicit gradient blocks give
\begin{equation}
 \|G\|\le\sqrt{\mathcal K},\qquad
 \|\dot G\|\le C_gB_n\rho,\qquad
 C_g=J_*\left(1+\sum_{j=2}^LH_{j-1}\right)
    +2S^2\sum_{j=2}^L\tau_jq_{j-1}^*+2Sq_L^*.
 \label{cp:src-gradient-derivative}
\end{equation}
For the second estimate differentiate each column block
\(m^{-1/2}\delta_a^{(1)}v_a^\top/\sqrt n\),
\(m^{-1/2}\delta_a^{(j)}h_a^{(j-1)\top}/n\), and
\(m^{-1/2}h_a^{(L)}/\sqrt n\), then bound Frobenius norm.
The division by \(\sqrt m\) prevents a sample factor.

The exact residual equation is \(\dot r=-2G^\top Gr\).
On a vertical segment based at real \(t\ge0\), norm
differentiation and real fitting give
\begin{equation}
 \rho(t\pm iv)\le Ye^{-\nu t}e^{2\mathcal Kv},\qquad
 \int_0^{h(t)}\rho(t\pm iv)\,dv\le Yr_n.
 \label{cp:src-vertical}
\end{equation}
The second bound is the defining identity for \(h(t)\).  It also
gives \(\rho\le Y+2\mathcal KYr_n\le2Y\) eventually.
The base variational equation is \(\dot P=-2GG^\top P\), with
algebraic transposes.  At a real anchor,
\(A_t=2G(t)G(t)^\top\) is real symmetric, so \(-iA_t\)
generates unitary motion.  Freeze this generator on the vertical
segment and use \eqref{cp:src-gradient-derivative}:
\[
 \|2G(t\pm iv)G(t\pm iv)^\top-A_t\|
 \le4\sqrt{\mathcal K}C_gB_n\int_0^v\rho(t\pm iu)\,du.
\]
Conjugating by the frozen unitary bounds its propagator by
\[
 \exp\left(4\sqrt{\mathcal K}C_gB_n
       \int_0^{h(t)}(h(t)-v)\rho(t\pm iv)\,dv\right)
 \le\exp\left(\frac{2C_g}{\sqrt{\mathcal K}}B_nYr_n\right).
\]
For the last inequality, integrate the exponential bound in
\eqref{cp:src-vertical}; the double integral is at most
\(Yr_n/(2\mathcal K)\).  Products on ordered disjoint subsegments
obey the same bound, since their growth integrals add.  Forward real
segments contract.  The negative-real corner has total exceptional
length at most \(2r_n\), costing at most \(e^{4\mathcal Kr_n}\).
Hence impose the directly checkable eventual gate
\begin{equation}
 4\mathcal Kr_n+(2C_g/\sqrt{\mathcal K})B_nYr_n<\log(3/2).
 \label{cp:src-flaring-gate}
\end{equation}
It supplies the factor-two base premise, without treating a complex
Gram matrix as positive.  Its left side tends to zero because
\(B_nr_n=O(\log(e+\ell)/\sqrt\ell)\).

Also impose \eqref{cp:src-rectangle-gates} with
\(c_t=a/(64YSU_{\rm fin})\).  Additional activity is at most
\(8Yr_n\), and normalized parameter increments are at most
\(8YSD_Wr_n\le1/4\).  Thus operators, activity and
\(\rho/S\le\lambda/8\) meet the shared premises.
All canonical anchor contours have length \(O(\ell)\), and the
bounding rectangle
\([-h(0),T_+]+i[-h(T_+),h(T_+)]\) has sides \(O(\ell)\).
The local map/remainder argument is therefore unchanged: the
one-dimensional control entropy is
\(n^{5/8}\operatorname{poly}(\ell)\), endpoint meshes are
polynomial, and the same \(n^{.78}\) Gaussian tail dominates them.
After the improved response bound, vertical displacement is at most
\(2YSU_{\rm fin}\sqrt\ell r_n=a/32\); the negative corner costs
at most \(a/8\).  The pole and passive-size stops improve, and
cavity transfer is strict.

Finally Claim~\ref{cp:src-time-derivative}, integrated over these
segments, gives the complete-path normalized correction radius
\[
 D_n=4Yr_n\max\{2S\max_jq_j^*,J_*B_n/S\}
                 =O(\log(e+\ell)/\sqrt\ell).
\]
The commuting retractions preserve this estimate and give global
modulus \(O(B_n)\).  Thus the shared Gaussian moment and budget
argument applies.  Local holomorphic continuation glues because the
domain is vertically fibered, retracts onto its real interval, and
is simply connected; strict stopped margins continue across its
compact boundary.  This proves the source event and its confidence.

For the disk claim, if its real coordinate is nonnegative, then
\(h(x)\ge h(t)-d_h\varrho(t)\ge\varrho(t)\).
If it is negative, \(t<\varrho(t)\) and
\(h(t)\le h(0)+d_ht\) imply \(\varrho(t)<h(0)\).
The right endpoint is below \(T+h(T)<T_+\).
The radius is nondecreasing with Lipschitz constant
\(C_r=d_h/[2(1+d_h)]\).  Each complete adaptive panel contributes
at least \(1/(2+C_r)\) to \(\int dt/\varrho(t)\), by bounding
the radius on that panel by \((1+C_r/2)\varrho(t_b)\).
Consequently \(J\le1+5\int_0^Tdt/h(t)\).
Split at \(2\mathcal Kr_ne^{\nu t}=1\).  Before this point use
\(\log(1+x)\ge x/2\); afterward the logarithm grows at least
with slope \(\nu/2\).  This bounds the integral by
\[
 \frac2{\nu r_n}+\frac{4\mathcal K}{\nu}
             \log\left(1+\frac{\nu T}{2\log2}\right).
\]
Substitute \(\nu=\lambda/4\), \(S=16Y/\lambda\) and \(T\)
to obtain \eqref{cp:src-panel-count}.  The coefficient bounds follow
from Claim~\ref{cp:src-powers} and \(a^{-1}\le\beta/16\).
\end{proof}

\paragraph{5. Reduction to the public bounds.}
We now compare the stronger working estimates with the statement, so no
coefficient from this proof is needed by a later construction.  The working
whole-sphere time radius is
\[
 \frac{a}{64YSU\sqrt\ell}
   =\frac{a\lambda}{1024Y^2U\sqrt\ell}.
\]
Since \(U\le\beta^{26L}\sqrt{d+3}\) and \(a^{-1}\le\beta/16\),
\[
 \frac{1024U}{a}
 \le64\beta^{26L+1}\sqrt{d+3}
 \le\beta^{30L}\sqrt{d+3}.
\]
The last inequality uses \(64\le\beta^{4L-1}\), valid for
\(L\ge2\), \(\beta\ge10\).  This gives at least the public time
radius.  Likewise \(V\le\beta^{2L}\sqrt{d+3}\) gives
\[
 \max\{8,8V/a\}
 \le\max\{8,\tfrac12\beta^{2L+1}\sqrt{d+3}\}
 \le\beta^{3L}\sqrt{d+3},
\]
so the working query tube contains the stated tube.

The local power bounds yield, for every layer,
\[
 H_j\le3\beta^{2j}\le\beta^{3L},\qquad
 \max(k_j,\tau_j)\le3\beta^{4L-1}\le\beta^{4L}.
\]
Multiplication by \(S=16Ym/\gamma\) gives the backward bounds.
The four-family coordinate envelope obeys
\[
 M_0=10\max_j(H_j,\tau_j)
 \le30\beta^{4L-1}\le\beta^{6L}.
\]
For all real times, the fitting statement bounds normalized operators by
nine and \(\|w\|_2/\sqrt n\le2Y/\sqrt\lambda\).  To check the
remaining population comparison directly, Minkowski and linear activation
growth give
\(\sqrt{Q^{(j)}_{11}}\le b+s\sqrt{Q^{(j-1)}_{11}}\), starting at
\(Q^{(0)}_{11}=1\).  The local recursion for \(H_j\) dominates this
recursion, so \(\lambda\le Q^{(L)}_{11}\le H_L^2\).
Hence \(2Y/\sqrt\lambda\le SH_L\).  The same forward and backward
recursions with operator bound nine give \(H_j,Sk_j,S\tau_j\)
on the whole real trajectory.  Their parameter limits give the endpoint
bounds as well.  This proves \eqref{cp:source-rms} in both stated domains.

The working finite-query claim has exactly the public time radius.  Its
prediction envelope is \(SH_L^2\le16\beta^{6L}Ym/\gamma\), proving
\eqref{cp:finite-time}.  Its all-time training-carrier envelope is
\(2SK_{\rm src}\sqrt\ell\le32\beta^{21L}(Ym/\gamma)\sqrt\ell\).
Substituting the explicit definition of \(k^{(j)}\) proves
\eqref{cp:carrier}, without exporting that auxiliary notation.

For the finite panel, choose \(r_b=\varrho(t_b)\) in the working
domain's adaptive partition.  Its disk containment and half-radius
intervals are exactly those in (iii), and its source envelopes have just
been bounded in canonical quantities.  Its defining coefficients depend
only on the width, \(m,\gamma,Y\), activations and depth, so these choices
are deterministic.  Finally
\[
 40960\frac{U_{\rm fin}}a\le40960\beta^{27L}\le\beta^{30L},
 \qquad 80\mathcal K\le80\beta^{14L}\le\beta^{16L},
\]
because \(40960\le\beta^{3L}\) and \(80\le\beta^{2L}\).
Applying these inequalities to \eqref{cp:src-panel-count} proves
\eqref{cp:panel-source}.  For any requested finite collection of the
source domains, split the fixed failure allowance and intersect their
events.  This supplies the single event asserted in the proposition and
completes its proof.
\end{proof}
